% !TEX program = xelatex
\documentclass[11pt,a4paper]{article}
\usepackage{xeCJK}
\usepackage{fontspec}
\setmainfont{Times New Roman}
\setCJKmainfont{SimSun}
\usepackage{amsmath,amssymb}
\usepackage{graphicx}
\usepackage{geometry}
\usepackage{hyperref}
\geometry{margin=1in}
\title{遥感目标检测 — 高级基础知识与流程（PDF 版）}
\date{\today}
\begin{document}
\maketitle

\section*{说明}
本文档以中文详细说明：开放提示词、数据集规模与 GPU 估计、上/下采样、IoU/GIoU 用法、NMS、伪标签置信度与一致性方法、batch 含义与影响，以及从输入到输出的完整流程并给出数值示例。公式均使用 LaTeX 书写以确保可读与可编译。

\section{什么是开放提示词（Open Prompt / Prompt）}
开放提示词（Prompt）最初来自语言模型，是一段用于引导模型产生特定行为的“提示”。在遥感检测中，开放提示词可以扩展为：类别描述、场景元信息、可学习的条件向量或生成器的控制向量。

数学上，如果模型为 $f(x;\theta,p)$，其中 $x$ 为图像，$\theta$ 为模型权重，$p$ 为 prompt，则不同 prompt 会得到不同条件预测：
\[
y_p = f(x;\theta,p).
\]

当我们拥有一个 prompt 集合 $P=\{p_1,\dots,p_N\}$（即 prompt-bank）时，可按不同策略使用它：平均化、加权或检索最优 prompt。详见后文 "Prompt-bank" 部分。

\section{数据集大小与预计 GPU 耗费}
\subsection*{数据集举例}
- DOTA v2（近似规模）：约一万级原始图像；通过瓦片化（例如每图切 1024x1024 瓦片）后，瓦片总数常扩增到数万至十万级（取决于 tile stride 与 overlap）。

\subsection*{GPU 内存估计（经验值）}
给定配置：Backbone=ResNet-50，输入 tile=1024$\times$1024，FPN，单阶段 Head。显存估计（近似）：
\begin{itemize}
  \item batch size per GPU = 1：需要约 8--10 GB；
  \item batch size per GPU = 2：需要约 12--16 GB；
  \item batch size per GPU = 4：需要约 20--28 GB。
\end{itemize}

估计原因：特征图分辨率大，FPN 与 detection head 会在多尺度上保存中间激活，导致显存占用迅速上升。若使用更轻的 backbone（MobileNet、EfficientNet-lite）或把 tile 缩小到 512，会显著降低显存需求。

如果使用多 GPU（例如 4 卡），可以把全局 batch size 分摊到每卡，从而在相同学习率策略下提高吞吐量。

\section{上采样与下采样对应什么}
在卷积神经网络中：
\begin{itemize}
  \item 下采样（downsampling）通常意味着空间分辨率下降，常见操作：步幅为 2 的卷积（stride=2）、池化（max/avg pool）。下采样增大感受野、降低计算与内存，但会损失高频/细节信息。
  \item 上采样（upsampling）为把低分辨率特征恢复到更高分辨率，常见操作：最近邻/双线性插值、反卷积（transposed convolution）、子像素重排（pixel shuffle）。FPN 中常用上采样把深层语义向浅层传递。
\end{itemize}

数学上，若 $F\in\mathbb{R}^{H\times W\times C}$ 为低分辨率特征，上采样 $U(\cdot)$ 会输出 $\tilde F=U(F)\in\mathbb{R}^{\alpha H\times \alpha W\times C'}$，其中 $\alpha$ 是放大倍数（通常 $\alpha=2$）。

\section{IoU 和 GIoU 在哪里用到}
\subsection*{IoU（Intersection over Union）}
IoU 用于：
\begin{itemize}
  \item 训练阶段衡量 anchor 与 GT 的匹配（Assignment）；
  \item 评估指标（检测 IoU 阈值如 AP@0.5）；
  \item 后处理（NMS 中计算两个预测框是否重复）。
\end{itemize}

IoU 定义为：
\[
\mathrm{IoU}(A,B)=\frac{\mathrm{area}(A\cap B)}{\mathrm{area}(A\cup B)}.
\]

\subsection*{GIoU（Generalized IoU）}
GIoU 作为 IoU 的改进，常用于目标回归损失以获得更稳定的梯度，定义：
\[
\mathrm{GIoU}(A,B)=\mathrm{IoU}(A,B)-\frac{\mathrm{area}(C)-\mathrm{area}(A\cup B)}{\mathrm{area}(C)},
\]
其中 $C$ 是包含 $A$ 和 $B$ 的最小闭包框。GIoU 常被用作回归损失的替代（或混合）以提升在非重叠框情形下的收敛性。

在代码中：IoU 用于 assignment 与 NMS，GIoU（或 CIoU/DIoU）用于回归损失项 $L_{box}$。

\section{NMS（Non-Maximum Suppression）详解}
NMS 的目的：在推理阶段去除对同一目标的重复检测。算法流程：
\begin{enumerate}
  \item 对预测框按置信度 $s_i$ 降序排序；
  \item 取最高分框 $b_1$ 放入保留集合；
  \item 从剩余盒子中删除与 $b_1$ 的 IoU $> T_{nms}$ 的盒子；
  \item 重复步骤 2--3 直到没有框。
\end{enumerate}

直观上，NMS 认为若两个框高度重合（IoU 很大），它们很可能是对同一目标的重复检测，只保留得分更高者。阈值 $T_{nms}$ 通常取 0.4--0.6。\par

Soft-NMS 替代硬删除为按 IoU 衰减分数：若框 $b_j$ 与当前保留框 $b_k$ 的 IoU 为 $o$, 则把 $s_j$ 更新为 $s_j \leftarrow s_j \cdot f(o)$，常用的衰减函数是高斯或线性函数；这样在密集目标场景下更稳健。

举例：两个框置信度分别 0.95 和 0.9，IoU=0.7，若 $T_{nms}=0.5$，则第二个框被删除；在 Soft-NMS 中第二个框的分数会衰减到例如 0.5，但仍可能被保留（取决于阈值）。

\section{伪标签置信度、阈值、模型一致性与时序教师（Mean Teacher）}
\subsection*{伪标签置信度如何计算}
对每个预测框，模型通常输出类别概率向量 $p=(p_1,\dots,p_C)$，置信度可取为：
\[
s=\max_k p_k.
\]
或取置信度为正类概率 $p_{\hat y}$，其中 $\hat y=\arg\max_k p_k$。对于多任务（回归+分类），也可用 $s= p_{\hat y}\cdot f_{box}(\text{box quality})$。

\subsection*{阈值设定}
阈值 $\tau$ 一般为经验设定（常见 0.5--0.95），更严格的阈值（如 0.9）能降低伪标签噪声但减少被采纳的样本数量。选择策略：
\begin{itemize}
  \item 固定阈值：简单直接；
  \item 动态阈值：按每类置信度分布自适应（例如每类取 top-K 或按分位数设阈）；
  \item 质量加权：同时考虑类别置信度与 bbox 质量（IoU/置信度乘积）。
\end{itemize}

若 $s\ge\tau$ 则该预测可作为伪标签（$\tilde y$）被加入训练集中，但通常还会为伪标签评分加权，即训练损失乘以权重 $w(s)$（例如 $w(s)=s$ 或常数小权重）。

\subsection*{模型一致性}
模型一致性指多个预测来源之间的一致程度：
\begin{itemize}
  \item 多模型一致性：用多个不同初始化/架构/检查点的模型预测同一无标签样本，只有当多数模型给出相似预测（位置与类别）时才接纳伪标签；
  \item 数据增强一致性：对同一图像做不同增强，若预测一致则说明可信度更高。
\end{itemize}

\subsection*{时序教师（Mean Teacher）}
Mean Teacher 方法维持一对模型：student 与 teacher。teacher 的权重 $\theta^{(t)}$ 通过 student 的历史权重指数移动平均（EMA）更新：
\[
\theta^{(t)}_{teacher} \leftarrow \alpha\,\theta^{(t-1)}_{teacher} + (1-\alpha)\,\theta^{(t)}_{student},
\]
其中 $\alpha\in[0,1)$（例如 0.99）。teacher 通常比 student 更稳定，作为伪标签生成器提供更可靠的目标预测。

\section{Batch 的含义与对硬件与性能的影响}
\subsection*{定义}
Batch（批次）是指一次前向/反向传播中同时处理的样本数量。若全局 batch size 为 $B$，每个 GPU 上的 batch size 为 $B_{gpu}=B/\#GPU$（在数据并行时）。

\subsection*{影响}
\begin{itemize}
  \item 显存需求：与 batch 大小成正比；\
  \item 梯度估计方差：较大 batch 可降低梯度噪声，训练更稳定；\
  \item 学习率关系：线性缩放规则（Linear Scaling Rule）建议当 batch 成倍增加时相应放大学习率：$\eta' = \eta \cdot (B'/B)$。但过大 batch 会导致泛化性下降（需调参）；\
  \item 吞吐量与速度：更大 batch 可提高 GPU 并行利用率，但受限于显存与 I/O。\
\end{itemize}

经验建议：在 1--8 的每卡 batch 范围内，通常可获得稳定训练与可控显存；对大尺度检测任务，每卡 batch 经常为 1 或 2。

\section{完整流程（从输入到输出）与具体数值示例}
下面按紧密逻辑顺序列出完整流程，并给出一个数值化的例子（便于理解）。

\subsection*{流程步骤}
\begin{enumerate}
  \item 数据准备：原始大图（例如 2000 张，每张 6000$\times$6000）按 tile=1024、stride=768 切片，得到约 $\sim$15k 瓦片，标注转换为瓦片坐标；
  \item 数据增强：随机水平/垂直翻转、随机尺度 (0.8--1.2)、颜色扰动；
  \item Backbone 提取特征：ResNet-50 输出层级特征 $C_3,C_4,C_5$；
  \item Neck（FPN）：融合得到多尺度特征 $P_3,P_4,P_5$；
  \item Head（Anchor-based）：在每个位置放置 anchors（例如 scales $\{32,64,128\}$ 和 ratios $\{0.5,1,2\}$），预测分类概率与回归偏移；
  \item Assignment：计算 anchor 与 GT 的 IoU，正负样本按 $T_{pos}=0.5$, $T_{neg}=0.4$ 分配；
  \item 损失计算与优化：分类用 Focal Loss ($\gamma=2,\alpha=0.25$)，回归用 Smooth L1 或 GIoU，优化器用 SGD（$\eta=0.01$, momentum=0.9）；
  \item 推理后处理：对所有预测应用阈值筛选（例如 $s\ge0.05$），再执行 NMS（$T_{nms}=0.5$）；
  \item 评估：计算 AP@50、AP@75 及 mAP；
  \item 若使用自训练：用当前 checkpoint 预测无标签数据，筛选 $s\ge\tau$（例如 $\tau=0.9$）的伪标签并加入训练或加权训练；重复训练。
\end{enumerate}

\subsection*{数值例子（端到端）}
设初始条件：2000 原图，tile 得到 15000 瓦片；使用 2 张 16GB GPU，ResNet-50，tile=1024，batch per GPU=2（全局 batch=4），训练 12 个 epoch。

训练期间的典型超参：\[\eta=0.01,\; momentum=0.9,\; weight\_decay=1e-4.\]

每次迭代计算：
\begin{itemize}
  \item 正样本数 $N_{pos}$ 在 batch 内波动，例如平均为 32；
  \item 分类损失为 Focal Loss，总体值随训练下降（示例：初始 1.8 -> 收敛 0.12）；
  \item 伪标签阶段：用 teacher 生成预测并以 $\tau=0.9$ 筛选，假设在 5000 张无标签瓦片上得到 12000 个高置信伪标签，按权重 $w(s)=s$ 加入训练。
\end{itemize}

最终输出：模型 checkpoint（权重文件），以及评估结果表（AP 各类与 mAP）。

\section*{结语与后续可做的工程映射}
如需，我可以把上述关键公式逐条注释到代码实现中（例如在 \texttt{src/openprompt\_rs/} 中），或把本 TeX 转为 Notebook（包含小例子、可运行的 pseudo-label 筛选代码片段与可视化）。

\vfill
\noindent 文件生成位置：\texttt{docs/setup/advanced_knowledge.tex}
\end{document}
